\chapter{Künftige Arbeiten}

\section{Zweck dieses Kapitels}

Das in dieser Arbeit vorgestellte AIM-Framework schafft eine mathematische
und begriffliche Grundlage für die Alignment-zentrierte Modellierung.

Viele Aspekte bleiben jedoch offen und bieten Möglichkeiten für weitere
Forschung, Entwicklung und praktische Umsetzung.

Dieses Kapitel umreißt mögliche Richtungen künftiger Arbeiten, von
mathematischen Erweiterungen bis zu Softwarearchitektur und
Ingenieuranwendungen.

\section{Erweiterung des dynamischen Modells}

\subsection{Von vereinfachter Dynamik zu Fahrzeugmodellen}

Die gegenwärtige Formulierung verwendet vereinfachte dynamische Größen wie
effektive Querbeschleunigung und Ruck.

Künftige Arbeiten sollten dies zu realistischeren Fahrzeugmodellen
erweitern, darunter:
\begin{itemize}
\item Mehrkörper-Fahrzeugdynamik,
\item Rad--Schiene-Interaktion,
\item Federungsverhalten,
\item Gleiselastizität.
\end{itemize}

\subsection{Geschwindigkeitsabhängige Modellierung}

Die Einbeziehung veränderlicher Geschwindigkeitsprofile
\[
v(s)
\]
eröffnet die Möglichkeit einer realistischeren Auswertung und Optimierung
von Alignments unter Betriebsbedingungen.

\section{Schwerpunkt-Trassierung}

\subsection{Begriffliche Entwicklung}

Die Idee, die Trajektorie eines dynamisch relevanten Punktes wie des
Schwerpunkts zu entwerfen, stellt eine grundlegende Erweiterung des
klassischen Alignment-Entwurfs dar.

Weitere Arbeiten sind erforderlich, um:
\begin{itemize}
\item die Abbildung von Gleisgeometrie auf Fahrzeugtrajektorie zu formalisieren,
\item vollständige dynamische Modelle einzubeziehen,
\item geeignete Zielfunktionen zu bestimmen.
\end{itemize}

\subsection{Integration in die Optimierung}

Eine wesentliche Forschungsrichtung ist die Kopplung der
Schwerpunkt-Trassierung mit Optimierungsmethoden. Sie ermöglicht
komfortorientierten Entwurf, dynamische Optimierung und die direkte
Auswertung der Fahrzeugantwort.

\section{Fortgeschrittene Optimierungsmethoden}

\subsection{SQP-Erweiterungen}

Während Sequential Quadratic Programming eine leistungsfähige Grundlage
bietet, können künftige Arbeiten zugeschnittene SQP-Formulierungen für
dünn besetzte Alignment-Strukturen, verbesserte Bedingungsbehandlung und
adaptive Gewichtungsstrategien umfassen.

\subsection{Globale Optimierung}

Probleme der Alignment-Optimierung sind häufig nicht konvex. Künftige
Forschung kann globale Optimierungsverfahren, Mehrstartstrategien und
hybride deterministisch--stochastische Methoden untersuchen.

\section{Modellierung auf Netzebene}

\subsection{Vom einzelnen Alignment zum Netz}

Reale Eisenbahnsysteme bestehen aus miteinander verbundenen Alignments,
die komplexe Netze bilden.

Künftige Arbeiten sollten das AIM-Framework auf Verzweigungsstrukturen,
Weichen und Kreuzungen sowie die korridorweite Optimierung erweitern.

\subsection{Topologische Bedingungen}

Die systematische Behandlung topologischer Bedingungen bleibt eine
wesentliche Herausforderung. Sie umfasst Knotenkonsistenz,
Mehrgleiskopplung und netzweite Machbarkeitsbedingungen.

\section{Bibliotheken für Bedingungen und Ziele}

\subsection{Formalisierung von Bedingungen}

Obwohl sich viele Eisenbahnregeln als Bedingungen ausdrücken lassen,
bleibt die vollständige Formalisierung regulatorischer Rahmenwerke eine
offene Aufgabe.

Künftige Arbeiten sollten die systematische Kodierung von Standards, die
Validierung an realen Projekten und die Parametrisierung für
unterschiedliche Eisenbahnsysteme anstreben.

\subsection{Entwurf von Zielfunktionen}

Der Entwurf von Zielfunktionen ist für die optimierungsbasierte
Alignment-Konstruktion entscheidend. Mögliche Erweiterungen umfassen
Kostenmodelle, Instandhaltungsoptimierung, Energieeffizienz und
Lebenszyklusbetrachtungen.

\section{Integration mit BIM und Datenstandards}

\subsection{Tiefere BIM-Integration}

Die Integration von AIM in BIM-Arbeitsabläufe kann durch engere Kopplung
mit IFC-Alignment-Modellen, bidirektionale Synchronisierung und die
Einbettung von AIM-Zuständen in BIM-Objekte weiterentwickelt werden.

\subsection{Erweiterungen von Standards}

Das AIM-Framework könnte zudem Erweiterungen bestehender Standards
anregen, beispielsweise die Aufnahme dynamischer Eigenschaften, die
Unterstützung von Optimierungsmetadaten und die Repräsentation
steuerungsbasierter Alignment-Modelle.

\section{Interaktive Systeme und Echtzeitsysteme}

\subsection{Live-Optimierung}

Eine besonders aussichtsreiche Richtung ist die Entwicklung von Systemen,
die während des Alignment-Entwurfs Rückmeldung in Echtzeit erlauben.
Dazu gehören Live-Vorschau während der Optimierung, interaktive Anpassung
von Bedingungen und unmittelbare Visualisierung dynamischer Wirkungen.

\subsection{Benutzerinteraktion}

Der Grabbeltisch-Begriff könnte auf interaktive Bedingungsdefinition,
visuelle Fehlersuche in Alignment-Zuständen und intuitive Bearbeitung von
Übergangsparametern erweitert werden.

\section{KI und assistierter Entwurf}

\subsection{Datengetriebene Ansätze}

Verfahren des maschinellen Lernens könnten das AIM-Framework ergänzen,
indem sie aus bestehenden Alignments lernen, Anfangslösungen vorschlagen
und Muster in Entwurfsentscheidungen erkennen.

\subsection{Hybride Systeme}

Eine aussichtsreiche Richtung ist die Verbindung von physikbasierter
Modellierung (AIM), datengetriebenen Methoden (KI) und interaktiver
Benutzersteuerung.

\section{Wiederbelebung von AXTRAN und darüber hinaus}

Klassische Systeme wie AXTRAN haben die Machbarkeit eines
optimierungsbasierten Alignment-Entwurfs gezeigt.

Künftige Arbeiten könnten solche Systeme rekonstruieren und
verallgemeinern, auf Probleme der Netzebene erweitern und moderne
numerische sowie rechentechnische Verfahren integrieren.

\section{Softwarearchitektur und Einsatz}

\subsection{Skalierbarkeit}

Praktische Anwendungen erfordern skalierbare Implementierungen, die große
Datenmengen und komplexe Netze verarbeiten können.

\subsection{Mehransichtssysteme}

Künftige Systeme könnten synchronisierte Mehrfensteransichten,
kombinierte 2D/3D-Repräsentationen und die Integration von GIS- und
CAD-Funktionalität umfassen.

\section{Langfristige Perspektive}

Langfristig könnte das AIM-Framework zu einem neuen Paradigma der
Infrastrukturmodellierung beitragen, in dem das Alignment zur zentralen
Ordnungsstruktur wird, Geometrie, Dynamik und Optimierung vereinigt sind
und der Ingenieurentwurf zu einem interaktiven Berechnungsprozess wird.

\section{Zusammenfassung}
\begin{summary}
Das AIM-Framework eröffnet ein breites Spektrum von Forschungs- und
Entwicklungsrichtungen.

Dazu gehören Erweiterungen des dynamischen Modells, die BIM-Integration,
Optimierung auf Netzebene, Echtzeitsysteme und KI-assistierter Entwurf.

Gemeinsam weisen diese Richtungen auf eine Zukunft, in der das
Eisenbahn-Alignment nicht länger als statische Geometrie, sondern als
dynamisches, berechenbares und optimierbares System im Kern des
Infrastrukturwesens behandelt wird.
\end{summary}
